% ==============================================================================
% IGCSE 0478 PAST PAPER & WORKSHEET BLUEPRINT
% ==============================================================================
% INSTRUCTIONS FOR CLAUDE:
% 1. Use this structural skeleton for all generated worksheets, quizzes, and exams.
% 2. DO NOT modify the column widths of the longtable in the Mark Scheme.
% 3. DO NOT use commas or bare parentheses inside tcolorbox titles (e.g., stratbox).
% 4. All question sub-parts must use the [leftmargin=*, label=(\alph*)] format.
% 5. Append marks at the end of the question text formatted as \textit{[X marks]}.
% ==============================================================================

% IGCSE 0478 shared preamble -- structure only, NO topic content.
% \input this file as the FIRST line of every 05_output/*.tex (it includes
% \documentclass, so it must come before anything else). Defines the three
% custom environments every generated document must use exclusively.
% See CLAUDE.md "Strict Separation of Structure and Content".
\documentclass[11pt]{article}

\usepackage[a4paper, margin=2cm]{geometry}
\usepackage{xcolor}
\usepackage{tcolorbox}
\usepackage{enumitem}
\usepackage{amsmath}
\usepackage{hyperref}

\tcbuselibrary{skins, breakable}

% stratbox: an exam strategy / technique tip.
\newtcolorbox{stratbox}[1][Exam Strategy]{
  enhanced, breakable, colback=blue!5, colframe=blue!60!black,
  fonttitle=\bfseries, title=#1
}

% critbox: a common mistake / misconception warning.
\newtcolorbox{critbox}[1][Common Mistake]{
  enhanced, breakable, colback=red!5, colframe=red!60!black,
  fonttitle=\bfseries, title=#1
}

% scenario: an applied exam-style scenario or worked example.
\newtcolorbox{scenario}[1][Scenario]{
  enhanced, breakable, colback=green!5, colframe=green!40!black,
  fonttitle=\bfseries, title=#1
}

% Note: Ensure \usepackage{multirow} and \usepackage{longtable} are in the preamble

\begin{document}

\title{\bfseries <Document Title e.g., Week 4: Algorithm Design>\\ \large <Subtitle e.g., Topics 7 \& 10>}
\author{IGCSE Computer Science 0478}
\date{}
\maketitle

% ============================================================
% SECTION: IN-CLASS WORKSHEET / EXAM PAPER
% ============================================================
\section*{<Section Name e.g., In-Class Worksheet>}

\textbf{Time Allowed: <XX> Minutes \,|\, Total Marks: <XX>}

\vspace{1em}

% --- QUESTION BLOCK ---
\noindent\textbf{Question <Number>} \hfill \textit{Source: <Optional past paper code>}
\\
<Optional context paragraph describing the scenario. E.g., A submersible's logging software...>

\begin{enumerate}[leftmargin=*, label=(\alph*)]
  
  % Standard text question
  \item <Question text detailing the problem or task.> \textit{[<X> marks]}
  \vspace{<Provide logical vertical space for student answers, e.g., 2cm or 3cm>}

  % Question with pseudocode/code block
  \item <Question text introducing the code.> \textit{[<X> marks]}
  \begin{lstlisting}
<Insert Pseudocode or Code Here>
  \end{lstlisting}
  \vspace{2cm}

  % Nested sub-questions (Roman numerals)
  \item <Question text for a multi-part question.>
  \begin{enumerate}[label=(\roman*), leftmargin=*, noitemsep]
    \item <Sub-question text.> \textit{[<X> marks]}
    \vspace{1.5cm}
    \item <Sub-question text.> \textit{[<X> marks]}
    \vspace{1.5cm}
  \end{enumerate}

\end{enumerate}

\newpage

% ============================================================
% SECTION: MARK SCHEME
% ============================================================
\section*{Mark Scheme}

\subsection*{Mark Scheme: <Section Name e.g., In-Class Worksheet> (Question <Number>)}

% The table widths are strictly locked to perfectly span an A4 page with 2cm margins
\begin{longtable}{@{}p{1.4cm} p{13.8cm} p{1.2cm}@{}}
\toprule
\textbf{Q} & \textbf{Marking Point} & \textbf{Marks} \\
\midrule
\endhead

% --- MARK SCHEME ROW BLOCK ---
% Note for Claude: The first argument of \multirow is the number of marking points (rows) for this specific question.
\multirow{<Number of Marking Points>}{*}{<Question Label e.g., 1(a)>}
 & <First marking point criteria> & 1 \\
 & <Second marking point criteria> & 1 \\
 & <Third marking point criteria> & 1 \\
\midrule

% --- MARK SCHEME ROW BLOCK WITH CODE/LOGIC ---
\multirow{<Number of Marking Points>}{*}{<Question Label e.g., 1(b)>}
 & <First criteria, e.g. Correct running TOTAL values> & 1 \\
 & <Second criteria, e.g. Correct loop termination> & 1 \\
\midrule

% Continue adding blocks for each question part...

\bottomrule
\end{longtable}

\end{document}